% Source: http://tex.stackexchange.com/a/150903/23931
\documentclass[12pt]{article}
\usepackage[letterpaper,margin=1in]{geometry}
\usepackage{xcolor}
\usepackage{fancyhdr}
\usepackage{tgschola} % or any other font package you like

\pagestyle{fancy}
\fancyhf{}
\fancyhead[C]{%
  \footnotesize\sffamily
  \yourname\quad
  \textcolor{blue}{\youremail}}

\newcommand{\soptitle}{Personal Statement}
\newcommand{\yourname}{Zi Hang (Frank) Yin}
\newcommand{\youremail}{zi.h.yin@mail.mcgill.ca}

\newcommand{\statement}[1]{\par\medskip
  \underline{\textcolor{blue}{\textbf{#1:}}}\space
}

\usepackage[
  colorlinks,
  breaklinks,
  pdftitle={\yourname - \soptitle},
  pdfauthor={\yourname},
  unicode
]{hyperref}

\begin{document}

\begin{center}\LARGE\soptitle\\
\large of \yourname\ (Music Theory PhD applicant for Fall 2026)
\end{center}

\hrule
\vspace{1pt}
\hrule height 1pt
\bigskip

\noindent Dear Music Theory Area, 
\bigskip

I am writing to express my strong interest in the PhD program in Music Theory. My academic trajectory has been deliberately interdisciplinary, combining music theory, music production, computer science, and electronics. This background has shaped both how I ask questions about music and how I design methods to answer them. After exploring multiple fields at the various levels, I am confident that music theory is the intellectual center of my work and my first choice for doctoral study.

My training spans analytical, technological, and creative practices. I hold a master’s degree in music, completed alongside formal studies in computer science, which provided me with extensive experience in data-heavy research, corpus thinking, and computational analysis. In parallel, my background in audio has given me a technologically grounded approach to timbre, spectrum, and production-oriented musical parameters—areas that are increasingly central to contemporary music theory, particularly in the analysis of popular and recorded music.



Over the past several years, my academic path has not been straightforward. Rather than progressing through a single, neatly bounded field, I pursued multiple programs and commitments in parallel, a choice shaped as much by circumstance as by intention. This period required sustained effort, careful self-regulation, and a willingness to confront limitations directly; it was neither efficient nor easy, and I do not idealize it. At times, working across several domains exposed the costs of divided attention and the difficulty of maintaining depth under pressure. At the same time, this experience forced me to develop strategies for structuring my work, integrating methods across fields, and converting constraints into productive forms of inquiry. In retrospect, what initially appeared as a disadvantage became a source of methodological flexibility and critical perspective, allowing me to move between analytical, technical, and creative modes of thinking with greater awareness of their respective demands and limits.


Alongside my studies, I have held teaching roles at multiple educational levels, which situates my work within music theory’s established disciplinarity, in which pedagogy functions as a constitutive practice through which theoretical knowledge is reproduced, transmitted, and operationalized. In the same practice-facing spirit, I conducted field research in Europe on music and antisemitism, examining how discourse, institutions, and listening contexts condition musical meaning and reception in politically charged settings. {theatre say something; challenges in artistic fields}


I am grateful that, lately, I have been listening to myself and pursuing work on my own terms. Looking forward, I am willing to devote five to six years to the discipline of music theory, as passage of sustained focus on the topic that I proposed and a valuable period of training in which I can build depth, method, and a coherent scholarly voice. Within that period of training, I also want my work to remain attentive to the conditions under which musical knowledge is produced and circulated today: upheaval of AI, polarized public discourse around culture, and a persistent mismatch between academic training and professional pathways. I am motivated by the belief that music theory is well positioned to contribute to a broader paradigm shift in how music is analyzed, taught, and understood in the age of AI and large-scale data. I am particularly interested in developing theory-driven, bias-aware and context-sensitive, empirically grounded approaches that integrate computational methods while preserving interpretive depth.


While I remain active to apply to various sources of funding, I am open to the option of completely self-fund my studies.

I am eager to pursue doctoral study in an environment that values analytical depth, methodological experimentation, and engagement with contemporary musical practices. I believe my background in music theory, production, computation, and pedagogy prepares me to contribute meaningfully to your academic community and to the evolving discipline of music theory.



% \bigskip

Thank you kindly for your consideration; I look forward to hearing from you and to further discussing my candidacy.    

\bigskip

\noindent Respectfully,

\bigskip

\noindent Zi Hang (aka. Frank) Yin 

\end{document}